\documentclass[11pt]{article}

\usepackage[margin=1in]{geometry}
\usepackage{enumitem}
\usepackage{listings}
\usepackage{xcolor}
\usepackage{hyperref}

\lstdefinestyle{yalnixpseudo}{
  basicstyle=\ttfamily\small,
  keywordstyle=\color{blue!60!black},
  commentstyle=\color{gray!70!black},
  columns=fullflexible,
  keepspaces=true,
  frame=single,
  breaklines=true,
  showstringspaces=false
}

\title{Yalnix Checkpoints 1 and 2 Pseudocode}
\author{COSC 58}
\date{\today}

\begin{document}
\maketitle

\section{Checkpoint 1 Scope}

Checkpoint 1 is a design checkpoint. The implementation should have real source files,
data structure sketches, function stubs, and pseudocode for the traps, syscalls, and
major kernel flows. The first version of these notes is reflected in the project source
tree and summarized here for Overleaf.

\subsection{Primary Source Files}

\begin{itemize}[noitemsep]
  \item \texttt{kernelstart.c}: boot sequence and transition into idle.
  \item \texttt{memory.c}, \texttt{memory.h}: physical frame tracking and page-table helpers.
  \item \texttt{process.c}, \texttt{process.h}: PCB layout and process ownership.
  \item \texttt{trap.c}, \texttt{trap.h}: interrupt vector setup and trap handlers.
  \item \texttt{syscalls.c}, \texttt{syscalls.h}: syscall dispatch and early syscall stubs.
  \item \texttt{idle.c}, \texttt{idle.h}: idle loop.
  \item \texttt{Queue.c}, \texttt{Queue.h}: intrusive queue helpers for ready and blocked queues.
  \item \texttt{kernelbrk.c}: kernel heap break management.
  \item \texttt{template.c}: future \texttt{LoadProgram} implementation plan.
\end{itemize}

\section{Core Data Structures}

\subsection{PCB}

Each process is represented by a process control block.

\begin{lstlisting}[style=yalnixpseudo]
PCB:
  pid
  state
  region1_pt
  saved UserContext
  saved KernelContext
  kernel_stack_pfns
  parent
  children queue
  queue entry for ready/blocked queues
\end{lstlisting}

\subsection{Global Kernel State}

\begin{lstlisting}[style=yalnixpseudo]
global current_process
global idle_process
global ready_queue
global kernel_region0_pt
global frame_allocated[]
global physical_frame_count
global vm_enabled
\end{lstlisting}

\section{Checkpoint 2 Boot Flow}

\subsection{\texttt{KernelStart}}

\begin{lstlisting}[style=yalnixpseudo]
KernelStart(cmd_args, pmem_size, uctxt):
  trace boot start

  ProcessInit()
    initialize ready queue

  MemoryInit(pmem_size)
    compute number of physical frames
    mark all frames free
    reserve kernel text/data/original heap frames
    allocate Region 0 page table

  MemoryBuildKernelRegion0()
    identity-map kernel text pages as read/exec
    identity-map kernel data and heap pages as read/write
    identity-map current kernel stack pages as read/write

  TrapInit()
    fill all trap vector entries with TrapUnhandled
    install TrapClock at TRAP_CLOCK
    install TrapKernel at TRAP_KERNEL
    write trap vector address to REG_VECTOR_BASE

  idle_process = ProcessCreateIdle(uctxt)
    allocate PCB
    allocate Region 1 page table
    allocate one frame for idle's user stack
    map top Region 1 page read/write
    ask helper_new_pid(region1_pt) for idle pid
    copy initial UserContext
    set pc = DoIdle
    set sp near VMEM_1_LIMIT

  current_process = idle_process

  MemoryInstallPageTables(idle_process->region1_pt)
    write REG_PTBR0, REG_PTLR0
    write REG_PTBR1, REG_PTLR1

  WriteRegister(REG_VM_ENABLE, 1)
  vm_enabled = true
  flush entire TLB

  copy idle_process->user_context into *uctxt
  return from KernelStart into user mode
\end{lstlisting}

\section{Virtual Memory}

\subsection{Initial Region 0}

Before enabling virtual memory, any physical page already used by the kernel must still
refer to the same memory after VM is enabled. The initial Region 0 page table therefore
uses identity mappings.

\begin{lstlisting}[style=yalnixpseudo]
MemoryBuildKernelRegion0():
  for page in [_first_kernel_text_page, _first_kernel_data_page):
    map page -> same pfn with PROT_READ | PROT_EXEC

  for page in [_first_kernel_data_page, _orig_kernel_brk_page):
    map page -> same pfn with PROT_READ | PROT_WRITE

  for page in kernel stack virtual pages:
    map page -> same pfn with PROT_READ | PROT_WRITE
\end{lstlisting}

\subsection{Frame Allocation}

\begin{lstlisting}[style=yalnixpseudo]
FrameAlloc():
  for each physical frame:
    if frame is free:
      mark allocated
      return pfn
  return ERROR

FrameFree(pfn):
  validate pfn
  mark frame free
\end{lstlisting}

\section{\texttt{SetKernelBrk}}

\begin{lstlisting}[style=yalnixpseudo]
SetKernelBrk(addr):
  requested_page = round addr up to page
  never shrink below _orig_kernel_brk_page

  if VM is not enabled:
    record highest requested kernel break
    return SUCCESS

  if growing:
    for each new heap page:
      allocate frame
      map Region 0 page read/write
      flush stale TLB entry

  if shrinking:
    for each removed heap page:
      free mapped frame
      invalidate PTE
      flush stale TLB entry

  update current kernel break
  return SUCCESS or ERROR
\end{lstlisting}

\section{Trap Handling}

\subsection{Trap Vector}

\begin{lstlisting}[style=yalnixpseudo]
TrapInit():
  for i in [0, TRAP_VECTOR_SIZE):
    trap_vector[i] = TrapUnhandled

  trap_vector[TRAP_CLOCK] = TrapClock
  trap_vector[TRAP_KERNEL] = TrapKernel

  WriteRegister(REG_VECTOR_BASE, trap_vector)
\end{lstlisting}

\subsection{Clock Trap}

\begin{lstlisting}[style=yalnixpseudo]
TrapClock(uctxt):
  save incoming UserContext into current PCB
  trace clock trap

  future checkpoint:
    update global clock tick
    wake delayed processes
    if ready queue is non-empty:
      context switch to next ready process
    else:
      keep or dispatch idle

  copy selected process UserContext back to *uctxt
\end{lstlisting}

\subsection{Kernel Trap}

\begin{lstlisting}[style=yalnixpseudo]
TrapKernel(uctxt):
  save incoming UserContext into current PCB
  trace syscall code in hex
  SyscallDispatch(uctxt)
  copy current PCB UserContext back to *uctxt
\end{lstlisting}

\subsection{Unhandled Trap}

\begin{lstlisting}[style=yalnixpseudo]
TrapUnhandled(uctxt):
  trace vector, code, fault address, pc, and sp
  future checkpoint:
    abort current user process if recoverable
    halt only for unrecoverable kernel state
\end{lstlisting}

\section{Early Syscalls}

\subsection{Dispatch}

\begin{lstlisting}[style=yalnixpseudo]
SyscallDispatch(uctxt):
  switch uctxt->code:
    YALNIX_GETPID:
      regs[0] = current_process->pid
    YALNIX_BRK:
      regs[0] = SysBrk(regs[0])
    YALNIX_DELAY:
      regs[0] = SysDelay(regs[0])
    default:
      trace unsupported syscall
      regs[0] = ERROR
\end{lstlisting}

\subsection{Future \texttt{Brk}}

\begin{lstlisting}[style=yalnixpseudo]
SysBrk(addr):
  validate requested break is in Region 1
  validate break does not collide with stack red zone
  if growing:
    allocate frames and map heap pages read/write
  if shrinking:
    free removed heap pages
  flush affected TLB entries
  return SUCCESS or ERROR
\end{lstlisting}

\subsection{Future \texttt{Delay}}

\begin{lstlisting}[style=yalnixpseudo]
SysDelay(clock_ticks):
  if clock_ticks < 0:
    return ERROR
  if clock_ticks == 0:
    return SUCCESS

  set process wake_tick = current_tick + clock_ticks
  move current process to delay queue
  context switch to next ready process or idle
  return SUCCESS when resumed
\end{lstlisting}

\section{Idle}

\begin{lstlisting}[style=yalnixpseudo]
DoIdle():
  forever:
    TracePrintf("DoIdle")
    Pause()
\end{lstlisting}

\section{Checkpoint 3 Notes}

The current \texttt{LoadProgram} file is intentionally a stub. The planned flow is:

\begin{lstlisting}[style=yalnixpseudo]
LoadProgram(name, args, proc):
  open executable
  use LoadInfo to validate Yalnix executable layout
  compute text, data, bss, heap, and stack pages
  copy arguments into temporary Region 0 buffer
  tear down old Region 1 mappings
  map new text pages read/write while loading
  map data/bss pages read/write
  map stack pages read/write at top of Region 1
  flush Region 1 TLB entries
  read text and data from executable
  zero bss
  copy argc/argv onto user stack
  change text protections to read/exec
  set proc->user_context.pc to entry point
  set proc->user_context.sp to new stack pointer
  return SUCCESS
\end{lstlisting}

\end{document}
